\begin{table}[H]
\centering
\scriptsize
\renewcommand{\arraystretch}{1.15}
\setlength{\tabcolsep}{5pt}
\begin{tabular}{|l|r|r|r|r|}
\hline
\multicolumn{5}{|c|}{\textbf{Gastos Financieros Servicio de Detección de Ofensividad}} \\ \hline
\textbf{TASA DE INTERÉS} & \multicolumn{4}{c|}{\textbf{12\%}} \\ \hline
\multicolumn{5}{|l|}{\textbf{CÁLCULO CON CUOTA NIVELADA}} \\ \hline
\textbf{AÑOS} & \textbf{INTERESES} & \textbf{CUOTA ANUAL} & \textbf{ABONO} & \textbf{SALDO} \\ \hline
Año 0 &  &  &  & 18,000.00 \\ \hline
Año 1 & 2,160.00 & 4,378.06 & 2,218.06 & 15,781.94 \\ \hline
Año 2 & 1,893.83 & 4,378.06 & 2,484.23 & 13,297.71 \\ \hline
Año 3 & 1,595.72 & 4,378.06 & 2,782.34 & 10,515.37 \\ \hline
Año 4 & 1,261.84 & 4,378.06 & 3,116.22 & 7,399.15 \\ \hline
Año 5 & 887.90 & 4,378.06 & 3,490.16 & 3,908.98 \\ \hline
Año 6 & 469.08 & 4,378.06 & 3,908.98 & 0.00 \\ \hline
\textbf{TOTAL INT.} & \textbf{S/ 8,268.38} &  &  &  \\ \hline
\end{tabular}

\vspace{0.2cm}
\textbf{Figura 10.13.} Gastos Financieros \textbf{[Elaboración Propia]}
\end{table}

\section{Costos Totales}

Los costos totales del servicio comprenden la suma de todos los egresos necesarios para el desarrollo, operación, administración, promoción y sostenimiento financiero del proyecto. Esta agrupación permite tener una visión integral del esfuerzo económico requerido para mantener en funcionamiento el servicio de detección de lenguaje ofensivo en redes sociales.

A continuación, en la Figura 10.14 se presenta un resumen de los principales componentes de costo anual, considerando la estructura de personal definida, los costos indirectos proyectados, los gastos administrativos, comerciales y financieros estimados a partir del primer año de operación.

\begin{table}[H]
\centering
\scriptsize
\renewcommand{\arraystretch}{1.15}
\setlength{\tabcolsep}{5pt}
\begin{tabular}{|l|r|r|r|r|r|r|}
\hline
\multicolumn{7}{|c|}{\textbf{COSTO TOTAL DEL SERVICIO}} \\ \hline
\textbf{Concepto} & \textbf{Año 1} & \textbf{Año 2} & \textbf{Año 3} & \textbf{Año 4} & \textbf{Año 5} & \textbf{Año 6} \\ \hline
Total Costo Desarrollo & 93,531 & 117,185 & 161,744 & 331,346 & 347,237 & 364,361 \\ \hline
Total Gasto Adm. & 85,509 & 89,288 & 118,321 & 186,020 & 195,020 & 204,704 \\ \hline
Total Gasto Ventas & 62,191 & 100,625 & 107,631 & 124,518 & 133,722 & 144,413 \\ \hline
\textbf{Costo de Venta} & \textbf{S/ 241,230.81} & \textbf{S/ 307,097.16} & \textbf{S/ 387,695.41} & \textbf{S/ 641,884.67} & \textbf{S/ 675,978.87} & \textbf{S/ 713,477.28} \\ \hline
\end{tabular}

\vspace{0.2cm}
\textbf{Figura 10.14.} Costo Total del Servicio \textbf{[Elaboración Propia]}
\end{table}

\section{Proyección a la inversión}

\subsection{Estimación de la Inversión Total}

Para el inicio de las operaciones se ha estimado hacer una inversión total de S/ 18,450.00 los cuales cubrirían la compra de los equipos para cada equipo del proyecto.

El equipo de desarrollo necesita equipos con alta GPU para poder entrenar modelos inteligentes. Debido a este requerimiento se entiende porque las computadoras personales resultan más costosas que la de los demás equipos.
